\chapter{\ifproject%
% \ifenglish Experimentation and Results\else การทดลองและผลลัพธ์\fi
% \else%
\ifenglish System Evaluation\else การประเมินระบบ\fi
\fi}

\section{การประเมินประสิทธิภาพของอุปกรณ์ไอโอที}

ในปัจจุบัน อุปกรณ์ไอโอทีได้รับความนิยมอย่างแพร่หลาย โดยเฉพาะในด้านการตรวจวัดสิ่งแวดล้อม เช่น การวัดอุณหภูมิ ความชื้นและคุณภาพอากาศเพื่อให้มั่นใจว่าอุปกรณ์สามารถทำงานได้อย่างมีประสิทธิภาพ การประเมินประสิทธิภาพของอุปกรณ์จึงเป็นสิ่งสำคัญ 
การศึกษานี้มีวัตถุประสงค์เพื่อประเมินประสิทธิภาพของเซนเซอร์วัดอุณหภูมิที่เชื่อมต่อกับ บอร์ด KidBright โดยพิจารณาจากปัจจัยสำคัญ ได้แก่ ความแม่นยำในการวัด ความเร็วในการตอบสนอง และความเสถียรของระบบ 
การทดสอบเซนเซอร์วัดอุณหภูมิที่เชื่อมต่อกับบอร์ด KidBright ในครั้งนี้ใช้วิธีการดังต่อไปนี้ 

\begin{enumerate}
    \item การทดสอบความแม่นยำ 
    \begin{itemize}
        \item ใช้เซนเซอร์วัดอุณหภูมิร่วมกับบอร์ด KidBright ในการเก็บค่าข้อมูลอุณหภูมิทุกๆ 1 วินาที 
    \end{itemize}
    \item การทดสอบความเร็วในการตอบสนอง 
    \begin{itemize}
        \item นำเซนเซอร์จากอุณหภูมิห้องไปยังแหล่งความเย็น (เช่น น้ำเย็น) และจับเวลาในการเปลี่ยนค่าของเซนเซอร์ 
        \item ทดสอบซ้ำหลายครั้ง 
    \end{itemize}
    \item การทดสอบความเสถียร 
    \begin{itemize}
        \item เก็บค่าข้อมูลเป็นระยะเวลา 4 ชั่วโมง และวิเคราะห์ว่าค่าที่ได้มีความผันผวนมากน้อยเพียงใด 
    \end{itemize}
    \item การทดสอบการใช้พลังงาน
    \begin{itemize}
        \item วัดกระแสไฟฟ้าที่เซนเซอร์และบอร์ด KidBright ใช้งาน โดยใช้เครื่องวัดพลังงานไฟฟ้า 
    \end{itemize}
\end{enumerate}

\section{การประเมินประสิทธิภาพของ Hyperledger Fabric} 

Hyperledger Fabric เป็นหนึ่งในแพลตฟอร์มบล็อกเชนที่ได้รับความนิยมในอุตสาหกรรม โดยมีคุณสมบัติสำคัญ เช่น ความสามารถในการควบคุมสิทธิ์การเข้าถึงข้อมูล ความปลอดภัย และการบันทึกข้อมูลแบบไม่สามารถแก้ไขได้ (Immutable Ledger) ซึ่งเหมาะสำหรับการใช้งานด้านไอโอทที 
การศึกษานี้มีวัตถุประสงค์เพื่อ ประเมินประสิทธิภาพของ Hyperledger Fabric ในการจัดเก็บข้อมูลอุณหภูมิจากเซนเซอร์ที่เชื่อมต่อกับบอร์ด KidBright ในแบบเรียลไทม์ โดยวัดประสิทธิภาพในด้าน ความเร็วในการเขียนข้อมูล, ความเร็วในการอ่านข้อมูล, ความสามารถในการรองรับธุรกรรมพร้อมกัน และความถูกต้องของข้อมูล 
การทดสอบ Hyperledger Fabric ในการจัดเก็บข้อมูลอุณหภูมิจากเซนเซอร์ที่เชื่อมต่อกับบอร์ด KidBright ในแบบเรียลไทม์ ในครั้งนี้ใช้วิธีการดังต่อไปนี้

\begin{enumerate}
    \item การตั้งค่าระบบทดสอบ 
    \begin{itemize}
        \item ใช้บอร์ด KidBright เชื่อมต่อกับเซนเซอร์วัดอุณหภูมิ (DS18B20) ส่งข้อมูลอุณหภูมิไปยัง Hyperledger Fabric ผ่าน MQTT หรือ HTTP API 
    \end{itemize}
    \item การทดสอบความเร็วในการเขียนข้อมูล 
    \begin{itemize}
        \item วัดระยะเวลาที่ใช้ตั้งแต่การส่งข้อมูลจาก KidBright จนถึงการบันทึกสำเร็จในบล็อกเชน 
    \end{itemize}
    \item การทดสอบความเร็วในการอ่านข้อมูล 
    \begin{itemize}
        \item ส่งคำขอดึงข้อมูลอุณหภูมิจาก Hyperledger Fabric และวัดระยะเวลาที่ใช้ 
    \end{itemize}
    \item การทดสอบความสามารถในการรองรับธุรกรรมพร้อมกัน (Throughput Test) 
    \begin{itemize}
        \item ใช้ Hyperledger Caliper ในการจำลองการส่งข้อมูลจากเซนเซอร์หลายตัว 
        \item วัดค่า TPS สูงสุดที่ระบบสามารถรองรับได้ 
    \end{itemize}
    \item การทดสอบความถูกต้องของข้อมูล 
    \begin{itemize}
        \item เปรียบเทียบค่าที่ได้รับจาก Hyperledger Fabric กับค่าจริงที่ส่งจากบอร์ด KidBright 
        \item ตรวจสอบว่ามีความคลาดเคลื่อนของข้อมูลหรือไม่ 
    \end{itemize}
\end{enumerate} 
